% Generated draft: requires mathematical and source review.
\documentclass[11pt]{article}
\usepackage[margin=25mm]{geometry}
\usepackage{fontspec}
\setmainfont{Latin Modern Roman}
\usepackage{amsmath,amssymb,mathtools,unicode-math}
\setmathfont{Latin Modern Math}
\usepackage{enumitem,needspace,xurl,hyperref}
\setlength{\parindent}{0pt}
\setlength{\parskip}{5pt}
\setlength{\emergencystretch}{3em}
\newcommand{\N}{\mathbb{N}}
\newcommand{\Z}{\mathbb{Z}}
\newcommand{\Q}{\mathbb{Q}}
\newcommand{\R}{\mathbb{R}}
\newcommand{\C}{\mathbb{C}}
\newcommand{\F}{\mathbb{F}}
\newcommand{\A}{\mathbb{A}}
\newcommand{\PP}{\mathbb{P}}
\newcommand{\E}{\mathbb{E}}
\newcommand{\eps}{\varepsilon}
\newcommand{\dd}{\,\mathrm{d}}
\begin{document}
\section*{UnsolvedMath ID: 11000069}
\textbf{Source problem code:} catalogue entry
\subsection{Definitions and mathematical statement}
For a genus-\(g\) surface, the usual level-\(L\) subgroup is
\[
 \operatorname{Mod}_g[L]=
 \ker\bigl(\operatorname{Mod}_g\to
             \operatorname{Aut}(H_1(\Sigma_g;\Z/L\Z))\bigr),
\]
with the image preserving the intersection form. The record does not state its genus, level, boundary or marking conventions; the displayed homological definition is the conventional interpretation requiring confirmation.

Problem 5.25 asks for an \emph{explicit finite presentation} of \(\operatorname{Mod}_g[L]\): a finite specified generating set and finite specified relators presenting that group. It does not supply such generators or relations, or define the required uniformity in \(g,L\). Abstract existence or a general finite-index rewriting procedure is not automatically the manageable explicit presentation sought for the following computations; the acceptance criterion must be recovered from context.

The source's motivation is a possible extension of Pitsch's method [Pi], which it says proves
\[
 \operatorname{rank}H_2(\operatorname{Mod}_{g,1};\Z)\le1
 \qquad(g\ge4).
\]
It reports an extension to punctured and bordered cases by Korkmaz--Stipsicsz, reference [Ko], preserving the supplied spelling. The exact meaning of the subscript \((g,1)\) and all theorem hypotheses are absent.

For a presentation \(G=F/R\), with \(F\) free and \(R\) the normal closure of the relators, the Hopf formula used is
\[
 H_2(G;\Z)\cong\frac{R\cap[F,F]}{[F,R]}.\tag{17}
\]
Thus one analyzes relators that lie in the commutator subgroup modulo commutators with relators. The source describes simplifying the numerator computation to the rank of an integer matrix and proposes a similar approach for the level subgroup, perhaps with computer computation, at least for small \(L\), seeking a general pattern. It does not assert that a presentation already yields a completed homology computation, a rank formula for level groups, or a uniform algorithmic complexity bound.
\subsection{Short English statement}
Find explicit finite generators and relators for the level subgroup, with an eye toward computing its second integral homology using Hopf's formula and integer matrices. The source gives no presentation or uniform pattern and only proposes the homology computation as a strategy.
\subsection{Sources}
\begin{itemize}
\item Local UnsolvedMath registry, permanent ID \(11000069\), exact source code \texttt{catalogue entry}, Problem 5.25, ``Presentation for level \(L\) structures,'' PDF page 52. No individual proposer or proposed year is supplied.
\item Farb, edited volume \emph{Problems on Mapping Class Groups and Related Topics}: \url{https://www.math.uchicago.edu/~farb/papers/mcgbook.pdf}. The volume has 329 explicitly typeset numbered items, with repeated local numbering distinguished by PDF page.
\item Supplied internal references: Pitsch [Pi] for the rank bound and presentation/Hopf method; Korkmaz--Stipsicsz [Ko] for punctured and bordered extensions. Complete bibliography entries are not included.
\item Extraction pdf-font-structure, accessed 2026-07-29. Source evidence NEEDS\_REVIEW presents an open problem and requires a current primary-literature audit. Supplied OPEN-TRIAGE summary reports presentations in special cases and says the general question remains open, without naming cases or further references.
\item Dataset attribution: Ulam AI / UnsolvedMath Contributors, CC BY 4.0. Source rights NEEDS\_REVIEW: publicly accessible author/publisher copy, redistribution terms requiring review.
\end{itemize}
\end{document}
